% GENERATED FILE. Edit canonical lessons, then run npm run generate:books.
% canonical-source-hash: fnv1a64:56ab362357b5879a
% canonical-lessons: GU-C198-patiyu, GU-C198-sucna, GU-C198-jaherat, GU-C198-manai, GU-C198-jokham

\chapter{Signs in the street}
\label{ch:gu-signs-in-the-street}
\hlchaptermodality{198}

\begin{chapteropening}
\textbf{By the end of this chapter:} \emph{I can read notices, adverts, bans and danger signs.}

Complete the last lesson of chapter 198: I can read notices, adverts, bans and danger signs.
\end{chapteropening}

\section[\texorpdfstring{\gu{પાટિયું} (pāṭiyũ)}{પાટિયું (pāṭiyũ)}]{\gu{પાટિયું} (pāṭiyũ) — signboard, plank}

\label{lesson:GU-C198-patiyu}



\begin{quote}
Before the new one: say the Gujarati for closed, then the Gujarati for to open up.
\end{quote}

\subsection*{You'll want to know: \gu{પાટિયું}}
\textbf{\gu{પાટિયું}} — \emph{pāṭiyũ} — \textquotedblleft{}signboard, plank\textquotedblright{}.

\begin{quote}
\textbf{\gu{દુકાનનું} \gu{પાટિયું} \gu{નવું} \gu{છે}.} — \emph{dukānnũ pāṭiyũ navũ che} — \textquotedblleft{}The shop's signboard is new.\textquotedblright{}
\end{quote}

\begin{cousinweb}
\textbf{\gu{પાટિયું}} (\emph{pāṭiyũ}) is a tadbhava. CDIAL lists it, glossed \textquotedblleft{}plank\textquotedblright{}, under Sanskrit \emph{paṭṭa}, \textquotedblleft{}slab, tablet\textquotedblright{} (CDIAL 7699).
\end{cousinweb}

\begin{grammarlens}[title={What you've built}]
One more word for everyday life at A2.
\end{grammarlens}

\subsection*{Your turn}
\begin{itemize}
\raggedright
  \item \emph{Say it:} \emph{pāṭiyũ}
  \item \emph{Say it:} \emph{pāṭiyũ}, once more
  \item \emph{Say it:} \emph{dukānnũ pāṭiyũ navũ che} — The shop's signboard is new.
  \item \emph{Recall:} say the Gujarati for closed, then the Gujarati for to open up, then say \emph{pāṭiyũ} again
\end{itemize}

\begin{tcolorbox}[breakable,colback=teal!4,colframe=teal!35!black,title={Before you move on}]
What does \gu{પાટિયું} mean? (\textquotedblleft{}Signboard, plank\textquotedblright{}.) Say it once more.
\end{tcolorbox}
\section[\texorpdfstring{\gu{સૂચના} (sūcnā)}{સૂચના (sūcnā)}]{\gu{સૂચના} (sūcnā) — announcement, notice}

\label{lesson:GU-C198-sucna}



\begin{quote}
Before the new one: say the Gujarati for to open up, then the Gujarati for signboard.
\end{quote}

\subsection*{You'll want to know: \gu{સૂચના}}
\textbf{\gu{સૂચના}} — \emph{sūcnā} — \textquotedblleft{}announcement, notice\textquotedblright{}.

\begin{quote}
\textbf{\gu{પાટિયા} \gu{પર} \gu{સૂચના} \gu{છે}.} — \emph{pāṭiyā par sūcnā che} — \textquotedblleft{}There is a notice on the board.\textquotedblright{}
\end{quote}

A feminine noun.

\begin{cousinweb}
\textbf{\gu{સૂચના}} (\emph{sūcnā}) is Sanskrit \emph{sūcanā}, \textquotedblleft{}indication, information\textquotedblright{}, taken over whole as a learned loan, a tatsama. CDIAL has no entry for it.
\end{cousinweb}

\begin{grammarlens}[title={What you've built}]
One more word for everyday life at A2.
\end{grammarlens}

\subsection*{Your turn}
\begin{itemize}
\raggedright
  \item \emph{Say it:} \emph{sūcnā}
  \item \emph{Say it:} \emph{sūcnā}, once more
  \item \emph{Say it:} \emph{pāṭiyā par sūcnā che} — There is a notice on the board.
  \item \emph{Recall:} say the Gujarati for to open up, then the Gujarati for signboard, then say \emph{sūcnā} again
\end{itemize}

\begin{tcolorbox}[breakable,colback=teal!4,colframe=teal!35!black,title={Before you move on}]
What does \gu{સૂચના} mean? (\textquotedblleft{}Announcement, notice\textquotedblright{}.) Say it once more.
\end{tcolorbox}
\section[\texorpdfstring{\gu{જાહેરાત} (jāherāt)}{જાહેરાત (jāherāt)}]{\gu{જાહેરાત} (jāherāt) — advertisement}

\label{lesson:GU-C198-jaherat}



\begin{quote}
Before the new one: say the Gujarati for signboard, then the Gujarati for announcement.
\end{quote}

\subsection*{You'll want to know: \gu{જાહેરાત}}
\textbf{\gu{જાહેરાત}} — \emph{jāherāt} — \textquotedblleft{}advertisement\textquotedblright{}.

\begin{quote}
\textbf{\gu{દીવાલ} \gu{પર} \gu{મોટી} \gu{જાહેરાત} \gu{છે}.} — \emph{dīvāl par moṭī jāherāt che} — \textquotedblleft{}There is a big advert on the wall.\textquotedblright{}
\end{quote}

Built on \textbf{\gu{જાહેર}}, \emph{jāher}, 'public'.

\begin{cousinweb}
\textbf{\gu{જાહેરાત}} (\emph{jāherāt}) is a loan built on Arabic \emph{ẓāhir}, \textquotedblleft{}manifest, public\textquotedblright{}, which reached Gujarati through Persian. Its Arabic root is \emph{ẓ-h-r}, \textquotedblleft{}to appear\textquotedblright{} (Lane).
\end{cousinweb}

\begin{grammarlens}[title={What you've built}]
One more word for everyday life at A2.
\end{grammarlens}

\subsection*{Your turn}
\begin{itemize}
\raggedright
  \item \emph{Say it:} \emph{jāherāt}
  \item \emph{Say it:} \emph{jāherāt}, once more
  \item \emph{Say it:} \emph{dīvāl par moṭī jāherāt che} — There is a big advert on the wall.
  \item \emph{Recall:} say the Gujarati for signboard, then the Gujarati for announcement, then say \emph{jāherāt} again
\end{itemize}

\begin{tcolorbox}[breakable,colback=teal!4,colframe=teal!35!black,title={Before you move on}]
What does \gu{જાહેરાત} mean? (\textquotedblleft{}Advertisement\textquotedblright{}.) Say it once more.
\end{tcolorbox}
\section[\texorpdfstring{\gu{મનાઈ} (manāī)}{મનાઈ (manāī)}]{\gu{મનાઈ} (manāī) — ban, prohibition}

\label{lesson:GU-C198-manai}



\begin{quote}
Before the new one: say the Gujarati for announcement, then the Gujarati for advertisement.
\end{quote}

\subsection*{You'll want to know: \gu{મનાઈ}}
\textbf{\gu{મનાઈ}} — \emph{manāī} — \textquotedblleft{}ban, prohibition\textquotedblright{}.

\begin{quote}
\textbf{\gu{મંદિરમાં} \gu{બૂટની} \gu{મનાઈ} \gu{છે}.} — \emph{mandirmā̃ būṭnī manāī che} — \textquotedblleft{}Shoes are not allowed in the temple.\textquotedblright{}
\end{quote}

A feminine noun.

\begin{cousinweb}
\textbf{\gu{મનાઈ}} (\emph{manāī}) is a loan that reached Gujarati from Arabic through Persian, but its Arabic source is not settled. It may rest on \emph{manʿ}, \textquotedblleft{}prevention, forbidding\textquotedblright{}, root \emph{m-n-ʿ}, or on the root \emph{n-h-y}, \textquotedblleft{}to forbid\textquotedblright{}; Lane gives both roots (Lane).
\end{cousinweb}

\begin{grammarlens}[title={What you've built}]
One more word for everyday life at A2.
\end{grammarlens}

\subsection*{Your turn}
\begin{itemize}
\raggedright
  \item \emph{Say it:} \emph{manāī}
  \item \emph{Say it:} \emph{manāī}, once more
  \item \emph{Say it:} \emph{mandirmā̃ būṭnī manāī che} — Shoes are not allowed in the temple.
  \item \emph{Recall:} say the Gujarati for announcement, then the Gujarati for advertisement, then say \emph{manāī} again
\end{itemize}

\begin{tcolorbox}[breakable,colback=teal!4,colframe=teal!35!black,title={Before you move on}]
What does \gu{મનાઈ} mean? (\textquotedblleft{}Ban, prohibition\textquotedblright{}.) Say it once more.
\end{tcolorbox}
\section[\texorpdfstring{\gu{જોખમ} (jokham)}{જોખમ (jokham)}]{\gu{જોખમ} (jokham) — danger, risk}

\label{lesson:GU-C198-jokham}



\begin{quote}
Before the new one: say the Gujarati for advertisement, then the Gujarati for ban.
\end{quote}

\subsection*{You'll want to know: \gu{જોખમ}}
\textbf{\gu{જોખમ}} — \emph{jokham} — \textquotedblleft{}danger, risk\textquotedblright{}.

\begin{quote}
\textbf{\gu{આ} \gu{રસ્તા} \gu{પર} \gu{જોખમ} \gu{છે}.} — \emph{ā rastā par jokham che} — \textquotedblleft{}This road is dangerous.\textquotedblright{}
\end{quote}

A neuter noun.

\begin{cousinweb}
\textbf{\gu{જોખમ}} (\emph{jokham}) is a tadbhava. CDIAL lists it, glossed \textquotedblleft{}risk, venture, responsibility, loss\textquotedblright{}, under Sanskrit \emph{yōgakṣēmá}, \textquotedblleft{}preservation of one's property; insurance charge\textquotedblright{} (CDIAL 10527).
\end{cousinweb}

\begin{grammarlens}[title={What you've built}]
One more word for everyday life at A2.
\end{grammarlens}

\subsection*{Your turn}
\begin{itemize}
\raggedright
  \item \emph{Say it:} \emph{jokham}
  \item \emph{Say it:} \emph{jokham}, once more
  \item \emph{Say it:} \emph{ā rastā par jokham che} — This road is dangerous.
  \item \emph{Recall:} say the Gujarati for advertisement, then the Gujarati for ban, then say \emph{jokham} again
\end{itemize}

\begin{tcolorbox}[breakable,colback=teal!4,colframe=teal!35!black,title={Before you move on}]
What does \gu{જોખમ} mean? (\textquotedblleft{}Danger, risk\textquotedblright{}.) Say it once more.
\end{tcolorbox}
